\documentclass[a4paper,twoside,12pt]{article}

%
%preprint HAL 
%
%cambiar en las rerferencias las minusculas por mayusculas
%

\usepackage[all]{xy}
\usepackage{xspace}
\usepackage{amsthm,tikz,hyperref,color,amsmath,graphicx,amsfonts,gensymb,mathdots,geometry,amssymb,tensor,enumitem,wasysym,makeidx,stmaryrd}
\usepackage[utf8]{inputenc}
\usepackage[T1]{fontenc}
\usepackage[english]{babel}
\usetikzlibrary{cd,arrows}

\def\Spec{ \operatorname{Spec}}
\def\Trace{ \operatorname{Trace}}
\def\GL{ \operatorname{GL}}
\def\SO{ \operatorname{SO}}
\def\O{ \operatorname{O}}
\def\F{\mathbb F}
\def\Fqbar{\overline{\F}_q}
\def\Ql{\overline{\mathbb Q}_{\ell}}
\def\C{\mathbb C}

\def\V{\mathbf V}
\def\G{\mathbf G}
\def\T{\mathbf T}
\def\B{\mathbf B}
\def\U{\mathbf U}
\def\P{\mathbf P}
\def\Q{\mathbf Q}
\def\L{\mathbf L}
\def\M{\mathbf M}
\def\X{\mathbf X}
\def\Y{\mathbf Y}

\def\Go{\mathbf G^{\circ}}
\def\To{\mathbf T^{\circ}}
\def\Bo{\mathbf B^{\circ}}
\def\Po{\mathbf P^{\circ}}
\def\Qo{\mathbf Q^{\circ}}
\def\Lo{\mathbf L^{\circ}}
\def\Mo{\mathbf M^{\circ}}

\def\GoF{\mathbf G^{\circ F}}
\def\ToF{\mathbf T^{\circ F}}
\def\BoF{\mathbf B^{\circ F}}
\def\PoF{\mathbf P^{\circ F}}
\def\QoF{\mathbf Q^{\circ F}}
\def\LoF{\mathbf L^{\circ F}}
\def\MoF{\mathbf M^{\circ F}}

%Puedes hacer una notación para los G^F normales, que se llame \GoF o algo así.

%\DeclareMathOperator{\Sym}{Sym}
\def\Hom{\operatorname{Hom}}
\def\RLPG{R_{\L\subset \P} ^{\G}}
\def\Res{\operatorname{Res}}
\def\Ind{\operatorname{Ind}}
\def\Irr{\operatorname{Irr}}
\def\Id{\operatorname{Id}}
\def\ad{\operatorname{ad}}

\newtheoremstyle{tdp}{}{}{\itshape}{}{\bfseries}{:}{.5em}{}
\newtheoremstyle{ddp}{}{}{}{}{\bfseries}{:}{.5em}{}
\newtheoremstyle{not}{}{}{}{}{\itshape}{:}{.5em}{}

\theoremstyle{tdp}
\newtheorem{definition}{Definition}[section]
\newtheorem{theorem}[definition]{Theorem}
\newtheorem{prop}[definition]{Proposition}
\newtheorem{corollary}[definition]{Corollary}
\newtheorem{lemma}[definition]{Lemma}

\theoremstyle{ddp}
\newtheorem{example}[definition]{Example}

\theoremstyle{not}
%\newtheorem*{note}{Note}
\newtheorem*{remark}{Remark}
%\newtheorem*{notation}{Notation}

\title{Summary thesis.}

\author{Sergio Cia}
%\date{25 octobre 2023}

%\makeindex

\begin{document}
 
\maketitle

In my thesis, I built an action of a Kac-Moody type Lie algebra over the category of unipotent representations of even-dimensional finite orthogonal groups, based on analogous constructions made by Dudas, Varagnolo and Vasserot for finite classical Lie groups (see \cite{DVV}).

The study of finite groups of Lie type is of capital importance on finite group theory, since they form one of the four families of finite simple groups: 
\begin{itemize}[label=$\bullet$]
	\item Cyclic groups of prime order;
	\item alternating groups of size $n\geq 5$;
	\item $26$ sporadic groups;
	\item \textbf{Finite groups of Lie type}. 
\end{itemize}

Among the finite groups of Lie type, we find the finite classical groups, that mirror the structure of real Lie groups but over a finite field. In this subfamily we find, for example, the general linear (type $A_n$) and orthogonal groups (type $B_n$, $D_n$ and $\tensor*[^2]{}{}D_n$). %reformular para que sea distinto al researc statement

During my thesis, I focused on the even orthogonal groups, since type $D_n$ is known to be problematic whenever we have to prove results for simple finite groups. Forn my work the main issue was that the finite even special orthogonal groups do not have a proper parametrization of irreducible unipotent representations.

In order to build the previously mentioned action over the category of unipotent modules, I had to adapt some known properties of representations of finite reductive groups -the finite groups of Lie type coming form connected reductive algebraic groups- to the finite orthogonal groups and/or more generally to grouos of Lie type coming from disconnected reductive groups. 

The first adaptation was a global Mackey formula for the Deligne-Lusztig induction and restriction operators. Digne and Michel (see \cite{DiMi}) defined these operators for disconnected reductive groups, and exhibed a Mackey formula that worked for Harish-Chandra like inductions and restrictions and for torus. Moreover, their formula was defined for left cosets of the identity component, rather than the whole group. 

The second adaptation was a combinatorial formula for the Lusztig induction to a finite orthogonal group from a twisted Levi subgroup. This formula was firstly stated for finite reductive groups of type $B_n$, $D_n$ and $\tensor*[^2]{}{}D_n$ by Asai, and properly proved by Malle (see \cite[\S 4.6]{GeMa} for a proof). %referenciar el artículo original?
Proving it required an extensive combinatorial analysis of the parameters of the unipotent representations of finite orthogonal groups.

%Up to other minor adaptations (as the Theorem on unitriangularity of decomposition matrices of unipotent blocks, proved by Brunat, Dudas and Taylor in \cite{BDT}), I managed to construct the categorical action, 

The categorical action constructed after these and other adaptations let us describe the crystal graph over the modular representations of finite orthogonal groups, allowing us to obtain an explicit structure of these representations and the branching rules -otherwise difficult to obtain. %- and its compatibility with Lie algebra actions by induction and restriction functors sheds light on the block structure and branching rules. %modular decomposition.

This action let us describe as well the endomorphism algebra of the induction of a weakly cuspidal (i.e. ''minimal'') irreducible representation as a Hecke algebra, with their parameters being obtained from the weight of the representation by the action of the Lie algebra. Finally, we can obtain some derived equivalences between blocks, opening the path to prove Broué's abelian defect group conjecture.

\bibliographystyle{plain}
\bibliography{phd_biblio}

\end{document}